\documentclass[12pt]{article}
\usepackage{amsmath}
\usepackage{amssymb}
\usepackage{geometry}
\geometry{a4paper, margin=1in}

\title{Homework 15.6 Triple Integrals}
\author{Tashfeen Omran}
\date{April 2026}

\begin{document}

\maketitle
\tableofcontents
\newpage

\section{Part 1: Comprehensive Introduction, Context, and Prerequisites}

\subsection{Core Concepts}
Triple integrals extend the concept of double integrals to three dimensions. While a double integral $\iint_D f(x,y) \, dA$ accumulates a quantity over a two-dimensional region $D$, a triple integral $\iiint_E f(x,y,z) \, dV$ accumulates a quantity over a solid, three-dimensional region $E$. 

The new differential variable introduced here is $dV$, which represents an infinitesimally small volume element. In rectangular coordinates, this volume element is the product of three infinitesimal lengths: $dV = dx \, dy \, dz$. To evaluate a triple integral, we apply Fubini's Theorem in three dimensions, expanding the integral into three nested iterated integrals. The most challenging conceptual leap in this topic is determining the limits of integration. The innermost limits are typically surfaces (functions of two variables), the middle limits are curves (functions of one variable), and the outermost limits are strictly constants.

\subsection{Intuition and Derivation}
Imagine a solid block of metal where the density is not uniform; perhaps it is heavier on one side than the other. If the density at any point $(x, y, z)$ is given by a function $\rho(x, y, z)$, how do we find the total mass of the block? 

We derive the formula by slicing the block into tiny, microscopic cubes, each with a volume $\Delta V = \Delta x \Delta y \Delta z$. Because the cubes are so small, the density inside a single cube is approximately constant. The mass of one tiny cube is simply its density multiplied by its volume: $\Delta m = \rho(x,y,z) \Delta V$. 

To find the total mass, we sum up the masses of all these tiny cubes. Taking the limit as the size of the cubes approaches zero, the Riemann sum perfectly transitions into a triple integral:
\[ \text{Total Mass} = \lim_{\Delta V \to 0} \sum \rho(x,y,z) \Delta V = \iiint_E \rho(x,y,z) \, dV \]
If the function being integrated is simply $f(x,y,z) = 1$, the integral sums up purely the geometric volume elements, resulting in the total volume of the region $E$.

\subsection{Historical Context and Motivation}
The foundational concepts of integral calculus were laid by Isaac Newton and Gottfried Wilhelm Leibniz in the 17th century. However, the rigorous formulation of multiple integrals was developed much later, prominently by Bernhard Riemann in the 19th century. 

The historical motivation for triple integrals was deeply rooted in classical mechanics and physics. Physicists needed precise mathematical tools to calculate the centers of mass, moments of inertia, and gravitational fields of continuous, solid, non-uniform bodies (like planets or engineered components). Simple geometry could not handle bodies with variable densities or complex, curved boundaries. Triple integrals were developed precisely to model these continuous three-dimensional summations.

\subsection{Key Formulas}
\begin{itemize}
    \item \textbf{General Triple Integral (Iterated):} 
    \[ \iiint_E f(x,y,z) \, dV = \int_{x=a}^{x=b} \int_{y=g_1(x)}^{y=g_2(x)} \int_{z=h_1(x,y)}^{z=h_2(x,y)} f(x,y,z) \, dz \, dy \, dx \]
    \item \textbf{Volume of a Solid Region $E$:} 
    \[ V = \iiint_E 1 \, dV \]
    \item \textbf{Cylindrical Volume Element (Modified for $x$-axis orientation):}
    If a surface has rotational symmetry about the $x$-axis, it is highly efficient to use $y = r\cos\theta$ and $z = r\sin\theta$. The volume element becomes:
    \[ dV = r \, dx \, dr \, d\theta \]
\end{itemize}

\subsection{Prerequisites}
To succeed in this chapter, you must be fully proficient in:
\begin{enumerate}
    \item \textbf{Double Integrals:} Setting up and evaluating integrals over Type I and Type II regions in 2D.
    \item \textbf{Integration Techniques:} Mastery of $u$-substitution, Integration by Parts, and evaluating integrals with algebraic limits.
    \item \textbf{3D Coordinate Geometry:} The ability to visualize and sketch 3D surfaces such as planes (e.g., $x+y+z=1$), paraboloids (e.g., $z = x^2+y^2$), and parabolic cylinders (e.g., $z = 1-x^2$).
\end{enumerate}

\newpage
\section{Part 2: Detailed Homework Solutions}

\subsection*{Problem 1}
\textbf{Evaluate the triple integral $\iiint_B xyz^2 \, dV$, where $B$ is the rectangular box $B = \{(x, y, z) \mid 0 \le x \le 1, -1 \le y \le 6, 0 \le z \le 4\}$, integrating first with respect to $y$, then $z$, and then $x$.}

\textbf{Solution:}
We are given explicit constant bounds for all three variables, meaning the region is a standard rectangular prism. The required order of integration is $dy \, dz \, dx$.
\[ \iiint_B xyz^2 \, dV = \int_{0}^{1} \int_{0}^{4} \int_{-1}^{6} xyz^2 \, dy \, dz \, dx \]
First, integrate with respect to $y$, treating $x$ and $z$ as constants:
\begin{align*} 
\int_{-1}^{6} xyz^2 \, dy &= xz^2 \left[ \frac{1}{2}y^2 \right]_{-1}^{6} \\
&= xz^2 \left( \frac{1}{2}(6)^2 - \frac{1}{2}(-1)^2 \right) \\
&= xz^2 \left( 18 - \frac{1}{2} \right) = \frac{35}{2} xz^2 
\end{align*}
Next, substitute this back and integrate with respect to $z$:
\begin{align*} 
\int_{0}^{4} \frac{35}{2} xz^2 \, dz &= \frac{35}{2} x \left[ \frac{1}{3}z^3 \right]_{0}^{4} \\
&= \frac{35}{2} x \left( \frac{64}{3} - 0 \right) = \frac{35 \times 64}{6} x = \frac{1120}{3} x 
\end{align*}
Finally, integrate with respect to $x$:
\begin{align*} 
\int_{0}^{1} \frac{1120}{3} x \, dx &= \frac{1120}{3} \left[ \frac{1}{2}x^2 \right]_{0}^{1} \\
&= \frac{1120}{6} (1 - 0) = \frac{560}{3} 
\end{align*}
\textbf{Final Answer:} $\frac{560}{3}$

\vspace{0.5cm}
\subsection*{Problem 2}
\textbf{Evaluate the following integral using three different orders of integration: $\iiint_E (xy + z^2) \, dV$, where $E = \{(x, y, z) \mid 0 \le x \le 6, 0 \le y \le 4, 0 \le z \le 5\}$.}

\textbf{Solution:}
Because $E$ is a rectangular box, all limits are constants, and Fubini's Theorem guarantees that any order of integration will yield the same result.

\textbf{Order 1: $dz \, dy \, dx$}
\[ \int_{0}^{6} \int_{0}^{4} \int_{0}^{5} (xy + z^2) \, dz \, dy \, dx \]
Inner Integral ($dz$):
\[ \left[ xyz + \frac{1}{3}z^3 \right]_{0}^{5} = 5xy + \frac{125}{3} \]
Middle Integral ($dy$):
\[ \int_{0}^{4} \left( 5xy + \frac{125}{3} \right) \, dy = \left[ \frac{5}{2}xy^2 + \frac{125}{3}y \right]_{0}^{4} = \frac{5}{2}x(16) + \frac{125}{3}(4) = 40x + \frac{500}{3} \]
Outer Integral ($dx$):
\[ \int_{0}^{6} \left( 40x + \frac{500}{3} \right) \, dx = \left[ 20x^2 + \frac{500}{3}x \right]_{0}^{6} = 20(36) + \frac{500}{3}(6) = 720 + 1000 = 1720 \]

\textbf{Order 2: $dy \, dz \, dx$}
\[ \int_{0}^{6} \int_{0}^{5} \int_{0}^{4} (xy + z^2) \, dy \, dz \, dx \]
Inner Integral ($dy$):
\[ \left[ \frac{1}{2}xy^2 + yz^2 \right]_{0}^{4} = \frac{1}{2}x(16) + 4z^2 = 8x + 4z^2 \]
Middle Integral ($dz$):
\[ \int_{0}^{5} (8x + 4z^2) \, dz = \left[ 8xz + \frac{4}{3}z^3 \right]_{0}^{5} = 40x + \frac{4}{3}(125) = 40x + \frac{500}{3} \]
Outer Integral ($dx$):
\[ \int_{0}^{6} \left( 40x + \frac{500}{3} \right) \, dx = \left[ 20x^2 + \frac{500}{3}x \right]_{0}^{6} = 720 + 1000 = 1720 \]

\textbf{Order 3: $dx \, dy \, dz$}
\[ \int_{0}^{5} \int_{0}^{4} \int_{0}^{6} (xy + z^2) \, dx \, dy \, dz \]
Inner Integral ($dx$):
\[ \left[ \frac{1}{2}x^2y + xz^2 \right]_{0}^{6} = \frac{1}{2}(36)y + 6z^2 = 18y + 6z^2 \]
Middle Integral ($dy$):
\[ \int_{0}^{4} (18y + 6z^2) \, dy = \left[ 9y^2 + 6yz^2 \right]_{0}^{4} = 9(16) + 6(4)z^2 = 144 + 24z^2 \]
Outer Integral ($dz$):
\[ \int_{0}^{5} (144 + 24z^2) \, dz = \left[ 144z + 8z^3 \right]_{0}^{5} = 144(5) + 8(125) = 720 + 1000 = 1720 \]
\textbf{Final Answer:} $1720$

\vspace{0.5cm}
\subsection*{Problem 3}
\textbf{Evaluate the iterated integral: $\int_{0}^{5} \int_{0}^{z^2} \int_{0}^{y-z} (2x - y) \, dx \, dy \, dz$.}

\textbf{Solution:}
Inner Integral (with respect to $x$):
\begin{align*} 
\int_{0}^{y-z} (2x - y) \, dx &= \left[ x^2 - yx \right]_{x=0}^{x=y-z} \\
&= (y-z)^2 - y(y-z) \\
&= (y^2 - 2yz + z^2) - (y^2 - yz) \\
&= -yz + z^2 = z^2 - yz
\end{align*}
Middle Integral (with respect to $y$):
\begin{align*} 
\int_{0}^{z^2} (z^2 - yz) \, dy &= \left[ z^2y - \frac{1}{2}y^2z \right]_{y=0}^{y=z^2} \\
&= \left( z^2(z^2) - \frac{1}{2}(z^2)^2z \right) - 0 \\
&= z^4 - \frac{1}{2}z^5
\end{align*}
Outer Integral (with respect to $z$):
\begin{align*} 
\int_{0}^{5} \left( z^4 - \frac{1}{2}z^5 \right) \, dz &= \left[ \frac{1}{5}z^5 - \frac{1}{12}z^6 \right]_{0}^{5} \\
&= \frac{5^5}{5} - \frac{5^6}{12} \\
&= 5^4 - \frac{15625}{12} = 625 - \frac{15625}{12} \\
&= \frac{7500 - 15625}{12} = -\frac{8125}{12}
\end{align*}
\textbf{Final Answer:} $-\frac{8125}{12}$

\vspace{0.5cm}
\subsection*{Problem 4}
\textbf{Evaluate the iterated integral: $\int_{0}^{1} \int_{y}^{2y} \int_{0}^{x+y} 18xy \, dz \, dx \, dy$.}

\textbf{Solution:}
Inner Integral (with respect to $z$):
\begin{align*} 
\int_{0}^{x+y} 18xy \, dz &= 18xy \left[ z \right]_{0}^{x+y} \\
&= 18xy(x+y) = 18x^2y + 18xy^2
\end{align*}
Middle Integral (with respect to $x$):
\begin{align*} 
\int_{y}^{2y} (18x^2y + 18xy^2) \, dx &= \left[ 6x^3y + 9x^2y^2 \right]_{x=y}^{x=2y} \\
&= \left( 6(2y)^3y + 9(2y)^2y^2 \right) - \left( 6(y)^3y + 9(y)^2y^2 \right) \\
&= \left( 6(8y^3)y + 9(4y^2)y^2 \right) - (6y^4 + 9y^4) \\
&= (48y^4 + 36y^4) - (15y^4) \\
&= 84y^4 - 15y^4 = 69y^4
\end{align*}
Outer Integral (with respect to $y$):
\begin{align*} 
\int_{0}^{1} 69y^4 \, dy &= \left[ \frac{69}{5}y^5 \right]_{0}^{1} = \frac{69}{5}
\end{align*}
\textbf{Final Answer:} $\frac{69}{5}$

\vspace{0.5cm}
\subsection*{Problem 5}
\textbf{Evaluate the iterated integral: $\int_{1}^{7} \int_{-1}^{2} \int_{-y}^{z} \frac{z}{y} \, dx \, dz \, dy$.}

\textbf{Solution:}
Inner Integral (with respect to $x$):
\begin{align*} 
\int_{-y}^{z} \frac{z}{y} \, dx &= \frac{z}{y} \left[ x \right]_{x=-y}^{x=z} \\
&= \frac{z}{y} (z - (-y)) = \frac{z}{y}(z + y) = \frac{z^2}{y} + z
\end{align*}
Middle Integral (with respect to $z$):
\begin{align*} 
\int_{-1}^{2} \left( \frac{z^2}{y} + z \right) \, dz &= \left[ \frac{z^3}{3y} + \frac{z^2}{2} \right]_{z=-1}^{z=2} \\
&= \left( \frac{8}{3y} + \frac{4}{2} \right) - \left( \frac{-1}{3y} + \frac{1}{2} \right) \\
&= \frac{8}{3y} + 2 + \frac{1}{3y} - \frac{1}{2} \\
&= \frac{9}{3y} + \frac{3}{2} = \frac{3}{y} + \frac{3}{2}
\end{align*}
Outer Integral (with respect to $y$):
\begin{align*} 
\int_{1}^{7} \left( \frac{3}{y} + \frac{3}{2} \right) \, dy &= \left[ 3 \ln|y| + \frac{3}{2}y \right]_{1}^{7} \\
&= \left( 3 \ln(7) + \frac{21}{2} \right) - \left( 3 \ln(1) + \frac{3}{2} \right) \\
&= 3 \ln 7 + \frac{21}{2} - 0 - \frac{3}{2} \\
&= 3 \ln 7 + \frac{18}{2} = 3 \ln 7 + 9
\end{align*}
\textbf{Final Answer:} $3 \ln 7 + 9$

\vspace{0.5cm}
\subsection*{Problem 6}
\textbf{Consider $f(x, y, z) = 6x$. (a) Express the triple integral $\iiint_E f(x, y, z) \, dV$ as an iterated integral for the solid region $E$ bounded by the parabolic cylinder $z = 1 - x^2$ and the plane $y + z = 2$ in the first octant context. (b) Evaluate the iterated integral.}

\textbf{Solution for (a):}
From the provided image, the solid $E$ sits on the $xy$-plane.
1. The bounds for $x$ run across the visible extent of the parabolic cylinder $z = 1-x^2$ where it intersects $z=0$, which is $x = -1$ to $x = 1$. The outer integral limits are printed as $-1$ to $1$.
2. The $z$ values are bounded below by the $xy$-plane ($z=0$) and above by the parabolic cylinder $z = 1 - x^2$.
3. The $y$ values are bounded below by the $xz$-plane ($y=0$) and above by the plane $y + z = 2$, which gives $y = 2 - z$.
Following the differentials order printed ($dy \, dz \, dx$), we write:
\[ \int_{-1}^{1} \int_{0}^{1-x^2} \int_{0}^{2-z} 6x \, dy \, dz \, dx \]
\textit{Blank answers for (a): Top outer is 1. Top middle is $1-x^2$. Top inner is $2-z$.}

\textbf{Solution for (b):}
Inner Integral ($dy$):
\[ \int_{0}^{2-z} 6x \, dy = 6x(2-z) = 12x - 6xz \]
Middle Integral ($dz$):
\begin{align*} 
\int_{0}^{1-x^2} (12x - 6xz) \, dz &= \left[ 12xz - 3xz^2 \right]_{0}^{1-x^2} \\
&= 12x(1-x^2) - 3x(1-x^2)^2 \\
&= 12x - 12x^3 - 3x(1 - 2x^2 + x^4) \\
&= 12x - 12x^3 - 3x + 6x^3 - 3x^5 \\
&= 9x - 6x^3 - 3x^5
\end{align*}
Outer Integral ($dx$):
\[ \int_{-1}^{1} (9x - 6x^3 - 3x^5) \, dx \]
Notice that the integrand is an odd function (every term has an odd power of $x$). Integrating an odd function over a symmetric interval $[-a, a]$ always results in $0$.
\textbf{Final Answer (b):} $0$

\vspace{0.5cm}
\subsection*{Problem 7}
\textbf{Consider $f(x, y, z) = 7$. (a) Express the triple integral over the solid $E$ bounded by the parabolic cylinder $z = 4 - y^2$ and the planes $z = 4 - x$ and $z = 4 + x$. (b) Evaluate the iterated integral.}

\textbf{Solution for (a):}
1. Looking at the parabolic cylinder $z = 4 - y^2$, its intersection with $z=0$ dictates the $y$ bounds. Setting $4-y^2 = 0 \Rightarrow y = \pm 2$. Thus, $y$ ranges from $-2$ to $2$.
2. For a fixed $y$, $z$ ranges from the floor $z=0$ up to the parabolic cylinder $z = 4 - y^2$.
3. The solid's front and back faces are formed by $z = 4 - x$ and $z = 4 + x$. Solving for $x$, we find the boundaries are $x = 4-z$ and $x = z-4$. Since the diagram points the front wedge to $x > 0$ and back to $x < 0$, the $x$ values go from $z-4$ to $4-z$.
The order is given as $dx \, dz \, dy$:
\[ \int_{-2}^{2} \int_{0}^{4-y^2} \int_{z-4}^{4-z} 7 \, dx \, dz \, dy \]
\textit{Blank answers for (a): Top middle is $4-y^2$. Top inner is $4-z$.}

\textbf{Solution for (b):}
Inner Integral ($dx$):
\[ \int_{z-4}^{4-z} 7 \, dx = 7[x]_{z-4}^{4-z} = 7(4-z - (z-4)) = 7(8-2z) = 56 - 14z \]
Middle Integral ($dz$):
\begin{align*} 
\int_{0}^{4-y^2} (56 - 14z) \, dz &= \left[ 56z - 7z^2 \right]_{0}^{4-y^2} \\
&= 56(4-y^2) - 7(4-y^2)^2 \\
&= 224 - 56y^2 - 7(16 - 8y^2 + y^4) \\
&= 224 - 56y^2 - 112 + 56y^2 - 7y^4 \\
&= 112 - 7y^4
\end{align*}
Outer Integral ($dy$):
Because $112 - 7y^4$ is an even function, we can simplify:
\begin{align*} 
\int_{-2}^{2} (112 - 7y^4) \, dy &= 2 \int_{0}^{2} (112 - 7y^4) \, dy \\
&= 2 \left[ 112y - \frac{7}{5}y^5 \right]_{0}^{2} \\
&= 2 \left( 224 - \frac{7(32)}{5} \right) = 2 \left( 224 - \frac{224}{5} \right) \\
&= 2 \left( 224 \times \frac{4}{5} \right) = \frac{1792}{5}
\end{align*}
\textbf{Final Answer (b):} $\frac{1792}{5}$

\vspace{0.5cm}
\subsection*{Problem 8}
\textbf{Evaluate $\iiint_E 7 \sin(y) \, dV$, where $E$ lies below $z = x$ and above the triangular region with vertices $(0,0,0)$, $(\pi,0,0)$, and $(0,\pi,0)$.}

\textbf{Solution:}
The region $E$ is bounded below by the $xy$-plane ($z=0$) and above by the plane $z=x$.
The base is a triangle in the $xy$-plane bounded by $x=0$, $y=0$, and the line connecting $(\pi,0)$ and $(0,\pi)$, which has the equation $x + y = \pi \Rightarrow y = \pi - x$.
We will set up the integral in the order $dz \, dy \, dx$.
\[ \int_{0}^{\pi} \int_{0}^{\pi-x} \int_{0}^{x} 7\sin(y) \, dz \, dy \, dx \]
Inner Integral ($dz$):
\[ \int_{0}^{x} 7\sin(y) \, dz = 7\sin(y) [z]_{0}^{x} = 7x\sin(y) \]
Middle Integral ($dy$):
\begin{align*} 
\int_{0}^{\pi-x} 7x\sin(y) \, dy &= 7x \left[ -\cos(y) \right]_{0}^{\pi-x} \\
&= 7x (-\cos(\pi-x) - (-\cos(0))) 
\end{align*}
Using the identity $\cos(\pi - x) = -\cos(x)$:
\[ = 7x (- (-\cos x) + 1) = 7x(\cos x + 1) = 7x\cos x + 7x \]
Outer Integral ($dx$):
\[ \int_{0}^{\pi} (7x\cos x + 7x) \, dx = \int_{0}^{\pi} 7x\cos x \, dx + \int_{0}^{\pi} 7x \, dx \]
For $\int_{0}^{\pi} 7x\cos x \, dx$, use Integration by Parts. Let $u=7x$ and $dv=\cos x \, dx$. Then $du=7 \, dx$ and $v=\sin x$.
\begin{align*} 
\int_{0}^{\pi} 7x\cos x \, dx &= \left[ 7x\sin x \right]_{0}^{\pi} - \int_{0}^{\pi} 7\sin x \, dx \\
&= (0 - 0) - \left[ -7\cos x \right]_{0}^{\pi} \\
&= 7\cos(\pi) - 7\cos(0) = 7(-1) - 7(1) = -14
\end{align*}
For the second part:
\[ \int_{0}^{\pi} 7x \, dx = \left[ \frac{7}{2}x^2 \right]_{0}^{\pi} = \frac{7\pi^2}{2} \]
Combining the parts:
\[ -14 + \frac{7\pi^2}{2} \]
\textbf{Final Answer:} $\frac{7\pi^2}{2} - 14$

\vspace{0.5cm}
\subsection*{Problem 9}
\textbf{Evaluate the triple integral $\iiint_E 4xy \, dV$, where $E$ lies under the plane $z = 1 + x + y$ and above the region in the $xy$-plane bounded by the curves $y = \sqrt{x}$, $y = 0$, and $x = 1$.}

\textbf{Solution:}
The region $E$ is bounded below by $z=0$ and above by $z = 1+x+y$. 
The base region $D$ in the $xy$-plane is bounded by $x=0$ to $x=1$, and for a given $x$, $y$ goes from $y=0$ to $y=\sqrt{x}$.
\[ \iiint_E 4xy \, dV = \int_{0}^{1} \int_{0}^{\sqrt{x}} \int_{0}^{1+x+y} 4xy \, dz \, dy \, dx \]
Inner Integral ($dz$):
\begin{align*} 
\int_{0}^{1+x+y} 4xy \, dz &= 4xy(1+x+y) = 4xy + 4x^2y + 4xy^2
\end{align*}
Middle Integral ($dy$):
\begin{align*} 
\int_{0}^{\sqrt{x}} (4xy + 4x^2y + 4xy^2) \, dy &= \left[ 2xy^2 + 2x^2y^2 + \frac{4}{3}xy^3 \right]_{y=0}^{y=\sqrt{x}} \\
&= 2x(\sqrt{x})^2 + 2x^2(\sqrt{x})^2 + \frac{4}{3}x(\sqrt{x})^3 \\
&= 2x(x) + 2x^2(x) + \frac{4}{3}x(x^{3/2}) \\
&= 2x^2 + 2x^3 + \frac{4}{3}x^{5/2}
\end{align*}
Outer Integral ($dx$):
\begin{align*} 
\int_{0}^{1} \left( 2x^2 + 2x^3 + \frac{4}{3}x^{5/2} \right) \, dx &= \left[ \frac{2}{3}x^3 + \frac{2}{4}x^4 + \frac{4}{3} \left( \frac{2}{7} \right) x^{7/2} \right]_{0}^{1} \\
&= \frac{2}{3} + \frac{1}{2} + \frac{8}{21}
\end{align*}
Find a common denominator, which is 42:
\[ \frac{28}{42} + \frac{21}{42} + \frac{16}{42} = \frac{65}{42} \]
\textbf{Final Answer:} $\frac{65}{42}$

\vspace{0.5cm}
\subsection*{Problem 10}
\textbf{Evaluate $\iiint_E 7x \, dV$, where $E$ is bounded by the paraboloid $x = 5y^2 + 5z^2$ and the plane $x = 5$.}

\textbf{Solution:}
The solid $E$ opens along the $x$-axis. The boundaries for $x$ are the paraboloid $5y^2 + 5z^2$ (lower bound) and the plane $x=5$ (upper bound).
The projection of this solid onto the $yz$-plane is found by setting the bounds equal: $5 = 5y^2 + 5z^2 \Rightarrow y^2 + z^2 = 1$. This is a circle of radius 1.
Because of the circular symmetry in the $yz$-plane, we use a modified cylindrical coordinate system where $y = r\cos\theta$ and $z = r\sin\theta$. 
In this system, $y^2+z^2 = r^2$. The boundaries for $x$ become $5r^2 \le x \le 5$. 
The volume element becomes $dV = dx \, (r \, dr \, d\theta)$.
\[ \int_{0}^{2\pi} \int_{0}^{1} \int_{5r^2}^{5} 7x \, dx \, r \, dr \, d\theta \]
Inner Integral ($dx$):
\begin{align*} 
\int_{5r^2}^{5} 7xr \, dx &= 7r \left[ \frac{1}{2}x^2 \right]_{5r^2}^{5} \\
&= \frac{7r}{2} \left( (5)^2 - (5r^2)^2 \right) = \frac{7r}{2} (25 - 25r^4) = \frac{175}{2} (r - r^5)
\end{align*}
Middle Integral ($dr$):
\begin{align*} 
\int_{0}^{1} \frac{175}{2} (r - r^5) \, dr &= \frac{175}{2} \left[ \frac{1}{2}r^2 - \frac{1}{6}r^6 \right]_{0}^{1} \\
&= \frac{175}{2} \left( \frac{1}{2} - \frac{1}{6} \right) = \frac{175}{2} \left( \frac{3}{6} - \frac{1}{6} \right) \\
&= \frac{175}{2} \left( \frac{2}{6} \right) = \frac{175}{6}
\end{align*}
Outer Integral ($d\theta$):
\[ \int_{0}^{2\pi} \frac{175}{6} \, d\theta = \frac{175}{6} [\theta]_{0}^{2\pi} = \frac{175}{6}(2\pi) = \frac{175\pi}{3} \]
\textbf{Final Answer:} $\frac{175\pi}{3}$

\vspace{0.5cm}
\subsection*{Problem 11}
\textbf{Evaluate the triple integral $\iiint_E x \, dV$, where $E$ is bounded by the paraboloid $x = 8y^2 + 8z^2$ and the plane $x = 8$.}

\textbf{Solution:}
This problem is structurally identical to Problem 10. We transition to cylindrical coordinates in the $yz$-plane ($y=r\cos\theta$, $z=r\sin\theta$). 
The intersection $8 = 8y^2+8z^2$ gives $r^2=1$, so $r$ ranges from 0 to 1.
$x$ ranges from the paraboloid $8r^2$ up to the plane $8$.
\[ \int_{0}^{2\pi} \int_{0}^{1} \int_{8r^2}^{8} x \cdot r \, dx \, dr \, d\theta \]
Inner Integral ($dx$):
\begin{align*} 
\int_{8r^2}^{8} xr \, dx &= r \left[ \frac{1}{2}x^2 \right]_{8r^2}^{8} = \frac{r}{2} (64 - 64r^4) = 32(r - r^5)
\end{align*}
Middle Integral ($dr$):
\begin{align*} 
\int_{0}^{1} 32(r - r^5) \, dr &= 32 \left[ \frac{1}{2}r^2 - \frac{1}{6}r^6 \right]_{0}^{1} \\
&= 32 \left( \frac{1}{3} \right) = \frac{32}{3}
\end{align*}
Outer Integral ($d\theta$):
\[ \int_{0}^{2\pi} \frac{32}{3} \, d\theta = \frac{32}{3} (2\pi) = \frac{64\pi}{3} \]
\textbf{Final Answer:} $\frac{64\pi}{3}$

\vspace{0.5cm}
\subsection*{Problem 12}
\textbf{Use a triple integral to find the volume of the given solid: the tetrahedron enclosed by the coordinate planes and the plane $7x + y + z = 3$.}

\textbf{Solution:}
To find the volume, we evaluate $\iiint_E 1 \, dV$. 
The solid is bounded by $x=0$, $y=0$, $z=0$, and $z = 3 - 7x - y$. 
The projection onto the $xy$-plane is found by setting $z=0$, yielding the line $7x + y = 3$, or $y = 3 - 7x$. 
Setting $y=0$ gives the $x$-intercept: $7x = 3 \Rightarrow x = \frac{3}{7}$.
Therefore, the boundaries are:
$0 \le x \le \frac{3}{7}$, $0 \le y \le 3-7x$, $0 \le z \le 3-7x-y$.
\[ V = \int_{0}^{3/7} \int_{0}^{3-7x} \int_{0}^{3-7x-y} 1 \, dz \, dy \, dx \]
Inner Integral ($dz$):
\[ \int_{0}^{3-7x-y} 1 \, dz = 3 - 7x - y \]
Middle Integral ($dy$):
\begin{align*} 
\int_{0}^{3-7x} (3 - 7x - y) \, dy &= \left[ (3-7x)y - \frac{1}{2}y^2 \right]_{0}^{3-7x} \\
&= (3-7x)(3-7x) - \frac{1}{2}(3-7x)^2 \\
&= (3-7x)^2 - \frac{1}{2}(3-7x)^2 = \frac{1}{2}(3-7x)^2
\end{align*}
Outer Integral ($dx$):
\[ \int_{0}^{3/7} \frac{1}{2}(3-7x)^2 \, dx \]
Let $u = 3-7x$, then $du = -7 \, dx$, and $dx = -\frac{du}{7}$. 
Limits: when $x=0, u=3$; when $x=3/7, u=0$.
\begin{align*} 
\int_{3}^{0} \frac{1}{2}u^2 \left( -\frac{du}{7} \right) &= \frac{1}{14} \int_{0}^{3} u^2 \, du \\
&= \frac{1}{14} \left[ \frac{1}{3}u^3 \right]_{0}^{3} \\
&= \frac{1}{14} \left( \frac{27}{3} \right) = \frac{9}{14}
\end{align*}
\textbf{Final Answer:} $\frac{9}{14}$


\newpage
\section{Part 3: In-Depth Analysis of Problems and Techniques}

\subsection{A) Problem Types and General Approach}
This homework set comprehensively spans four core categories of triple integral problems:
\begin{itemize}
    \item \textbf{Type 1: Setup and Fubini Iteration (Problems 1-5)}
    These questions present the integral entirely predefined with algebraic limits. The primary strategy here is rigid adherence to the order of integration. Work strictly from the innermost differential outward. You treat all outer variables as completely constant during the inner integrations.
    \item \textbf{Type 2: Graphic Region Translation (Problems 6, 7)}
    Here, a visual diagram is provided. The strategy is to mentally project the 3D solid onto a 2D coordinate plane to find the outer two limits. You identify the "floor" and "ceiling" functions to act as your inner integral bounds.
    \item \textbf{Type 3: Analytic Boundary Definitions (Problems 8, 9, 12)}
    These are the most rigorous Cartesian problems. You are given equations of planes or curves. The strategy is to systematically establish bounds by intersecting surfaces. Set $z=0$ to find the "shadow" region in the $xy$-plane, and treat that exactly like a double integral from Calculus II.
    \item \textbf{Type 4: Cylindrical Coordinate Conversions (Problems 10, 11)}
    Recognized instantly by equations combining squares of two variables, like $y^2+z^2$. The strategy is to abandon Cartesian coordinates to prevent integrating impossible square roots. Map the variables exhibiting circular symmetry into polar equations, remembering to multiply the integrand by the Jacobian determinant $r$.
\end{itemize}

\subsection{B) Key Algebraic and Calculus Manipulations}
\begin{itemize}
    \item \textbf{Exploiting Symmetries (Problem 6):}
    A powerful calculus trick is recognizing when an integral will evaluate to zero without doing the heavy lifting. In Problem 6, evaluating $\int_{-1}^{1} (9x - 6x^3 - 3x^5) \, dx$ was bypassed by recognizing that odd powers of $x$ (where $f(-x) = -f(x)$) integrated over a perfectly symmetrical domain $[-a, a]$ inherently cancel themselves out completely.
    \item \textbf{Integration by Parts (Problem 8):}
    The term $7x\cos x$ required the method of Integration by Parts ($\int u \, dv = uv - \int v \, du$). Selecting $u=7x$ ensures that the derivative $du=7 \, dx$ reduces the algebraic complexity, stripping away the polynomial term to leave a simple trigonometric integral.
    \item \textbf{Jacobian Determinant in Coordinate Shift (Problems 10, 11):}
    When converting to cylindrical coordinates $(y=r\cos\theta, z=r\sin\theta)$, the differential area element changes fundamentally. It is an absolute requirement to multiply by $r$, changing $dx \, dy \, dz$ to $dx \, r \, dr \, d\theta$. Failing to include this $r$ is the most common mathematical manipulation error in multivariable calculus.
    \item \textbf{$U$-Substitution for Polynomial Simplification (Problem 12):}
    Rather than agonizingly expanding the binomial $(3-7x)^2$ and integrating term by term, recognizing the linear inner function allowed for a clean $u$-substitution, compressing three lines of algebra into a single, elegant step.
\end{itemize}

\newpage
\section{Part 4: "Cheatsheet" and Tips for Success}

\subsection{Important Formulas Summary}
\begin{itemize}
    \item \textbf{Volume as a Triple Integral:} $V = \iiint_E 1 \, dV$
    \item \textbf{Average Value of a Function in 3D:} $f_{\text{avg}} = \frac{1}{V(E)} \iiint_E f(x,y,z) \, dV$
    \item \textbf{Cylindrical Coords (Standard):} $x = r\cos\theta, y = r\sin\theta, z = z \implies dV = r \, dz \, dr \, d\theta$
    \item \textbf{Mass from Density:} $m = \iiint_E \rho(x,y,z) \, dV$
\end{itemize}

\subsection{Tricks and Shortcuts}
\begin{itemize}
    \item \textbf{The Separable Integrand Cheat:} If your bounds are entirely constant (a rectangular box) AND your integrand can be factored into purely single-variable terms (like $xyz^2$ in Problem 1), you can split the triple integral into the product of three separate 1D single integrals:
    \[ \iiint f(x)g(y)h(z) \, dV = \left(\int f(x)dx\right)\left(\int g(y)dy\right)\left(\int h(z)dz\right) \]
    \item \textbf{The Symmetry Shortcut:} Always check the parity of your integrand over the bounds. If $f(-x, y, z) = -f(x, y, z)$ and the region $E$ is symmetric across the $yz$-plane ($x$ limits from $-a$ to $a$), the integral is instantly $0$.
\end{itemize}

\subsection{Common Pitfalls}
\begin{itemize}
    \item \textbf{Variables in the Outer Bounds:} The absolute golden rule of multiple integrals is that the \textit{outermost limits of integration must be constants}. If you end up with an $x$, $y$, or $z$ in your final answer, you have constructed the iterated integral bounds incorrectly.
    \item \textbf{The Phantom Differential:} Don't forget the $r$ when you shift to cylindrical coordinates. Write $dV = r \, dz \, dr \, d\theta$ immediately as a single unit block.
    \item \textbf{Confusing the order of limits:} When calculating a volume bounded between surfaces, say $z=x$ and $z=0$, always mentally verify which surface is physically "above" the other in the region of interest to avoid negative volumes.
\end{itemize}

\newpage
\section{Part 5: Conceptual Synthesis and The "Big Picture"}

\subsection{A) Thematic Connections}
The defining theme connecting Calculus I, II, and III is the "Accumulation of Infinitesimals." In Calc I, you accumulated linear slivers to find the area under a curve. In Calc II, you accumulated geometric 2D disks to find the volume of a solid of revolution. Now, in Calc III, you are using double integrals to accumulate infinitely thin 1D pillars to find a volume, and triple integrals to accumulate infinitely tiny 3D cubes to find a mass or hyper-volume. 
The mathematical engine—the Riemann Sum—is identical in all three cases. Only the dimensionality of the measuring unit (length vs. area vs. volume) changes.

\subsection{B) Forward and Backward Links}
\textbf{Backward Link:} Triple integrals are a strict, logical evolution of double integrals. Setting up the outer two limits of a triple integral is functionally identical to defining the bounds of a 2D base region in a double integral.

\textbf{Forward Link:} You are currently learning the \textit{grammar} required to speak the language of Vector Calculus. In the coming chapters, you will encounter the Divergence Theorem, which mathematically connects the flow of a vector field \textit{across} a 2D surface to the accumulation of "source" points \textit{inside} the 3D solid bounded by that surface. Evaluating the inside of that solid requires absolute mastery of triple integrals over complex regions.

\newpage
\section{Part 6: Real-World Application and Modeling}

\subsection{A) Concrete Scenarios in Financial Engineering and Statistics}
\textbf{Scenario 1: Joint Probability Density Functions in Risk Assessment}
In quantitative finance and actuarial science, the risks associated with multiple continuous variables (such as the simultaneous default risk of three interconnected corporations) are modeled using a joint probability density function $f(x,y,z)$. To find the probability that a catastrophic financial event occurs where all three variables fall within a specific critical region $E$, actuaries must evaluate the triple integral $\iiint_E f(x,y,z) \, dV$.

\textbf{Scenario 2: Pricing Basket Options}
A basket option is a financial derivative whose payout depends on the weighted sum of several underlying assets. If a basket contains three assets with prices $S_1, S_2, S_3$ modeled continuously, the expected payoff of the option requires integrating the payout function over the multi-dimensional probability space of the asset paths. This relies entirely on triple (and higher-dimensional) integration techniques.

\subsection{B) Model Problem Setup: Joint Density of Interest Rates and Inflation}
Let's model an economic risk scenario. An economist models three variables for a developing nation: 
$X$ = Exchange Rate Fluctuation, 
$Y$ = Inflation Rate, 
$Z$ = Central Bank Interest Rate. 

They determine the joint probability density function is:
\[ f(x,y,z) = C \cdot e^{-(x+y+z)} \]
for $x > 0, y > 0, z > 0$ (and $0$ otherwise), where $C$ is a normalization constant.
The economist needs to find the probability that the combined rate factors exceed a critical threshold $K$, which would trigger a sovereign debt crisis. The region $E$ where the crisis occurs is defined by $x + y + z > K$. 

The mathematical model required to find the probability $P$ of this crisis is:
\[ P(X+Y+Z > K) = \iiint_{x+y+z > K} C \cdot e^{-(x+y+z)} \, dx \, dy \, dz \]
To set this up as an iterated integral over the domain space (subtracting the non-crisis region from 1):
\[ P = 1 - \int_{0}^{K} \int_{0}^{K-x} \int_{0}^{K-x-y} C \cdot e^{-(x+y+z)} \, dz \, dy \, dx \]

\newpage
\section{Part 7: Common Variations and Untested Concepts}

The most significant concept missing from the provided homework set is the integration of solids using \textbf{Spherical Coordinates}. 

While the homework covered rectangular coordinates and touched implicitly on cylindrical (using polar bounds in a 2D plane), some solids are profoundly tedious to integrate unless converted to spherical coordinates $(\rho, \theta, \phi)$. These include true spheres, cones, and ice-cream-cone-shaped wedges.

\textbf{Spherical Coordinate Conversions:}
\[ x = \rho \sin\phi \cos\theta, \quad y = \rho \sin\phi \sin\theta, \quad z = \rho \cos\phi \]
\textbf{The Spherical Volume Element (Crucial):}
\[ dV = \rho^2 \sin\phi \, d\rho \, d\phi \, d\theta \]

\textbf{Worked Example:} Evaluate $\iiint_E e^{(x^2+y^2+z^2)^{3/2}} \, dV$, where $E$ is the unit ball $x^2+y^2+z^2 \le 1$.

\textit{Explanation:} Attempting this in Cartesian coordinates results in $\int_{-1}^{1} \int_{-\sqrt{1-x^2}}^{\sqrt{1-x^2}} \int_{-\sqrt{1-x^2-y^2}}^{\sqrt{1-x^2-y^2}} e^{(x^2+y^2+z^2)^{3/2}} \, dz \, dy \, dx$, which is essentially impossible to integrate directly.
By shifting to spherical coordinates, $x^2+y^2+z^2 = \rho^2$, so the integrand becomes $e^{(\rho^2)^{3/2}} = e^{\rho^3}$. The boundaries for a full sphere are $0 \le \rho \le 1$, $0 \le \phi \le \pi$, and $0 \le \theta \le 2\pi$.

\textit{Evaluation:}
\begin{align*} 
\iiint_E e^{(x^2+y^2+z^2)^{3/2}} \, dV &= \int_{0}^{2\pi} \int_{0}^{\pi} \int_{0}^{1} e^{\rho^3} \rho^2 \sin\phi \, d\rho \, d\phi \, d\theta 
\end{align*}
Because the bounds are constants and the integrand separates into factors of individual variables, we can split it:
\begin{align*}
&= \left( \int_{0}^{2\pi} 1 \, d\theta \right) \left( \int_{0}^{\pi} \sin\phi \, d\phi \right) \left( \int_{0}^{1} \rho^2 e^{\rho^3} \, d\rho \right) \\
&= (2\pi) \cdot [-\cos\phi]_{0}^{\pi} \cdot \left[ \frac{1}{3}e^{\rho^3} \right]_{0}^{1} \\
&= (2\pi) \cdot (1 - (-1)) \cdot \left( \frac{1}{3}e - \frac{1}{3} \right) \\
&= 2\pi(2)\left(\frac{e-1}{3}\right) = \frac{4\pi(e-1)}{3}
\end{align*}

\newpage
\section{Part 8: Advanced Diagnostic Testing: "Find the Flaw"}

Below are five solutions to triple integral problems. Each contains a subtle but critical mathematical flaw. Your task is to act as the grader: identify the error, explain it in one sentence, and solve it correctly.

\vspace{0.3cm}
\textbf{Problem 1: Fubini Limit Swapping}
Evaluate $\int_{0}^{1} \int_{0}^{x} \int_{0}^{y} 24x \, dz \, dy \, dx$.
\textit{Flawed Solution:} 
We can swap the order of integration to $dx \, dy \, dz$ to make the algebra easier. 
$\int_{0}^{1} \int_{0}^{x} \int_{0}^{y} 24x \, dx \, dy \, dz = \int_{0}^{1} \int_{0}^{x} \left[ 12x^2 \right]_{0}^{y} \, dy \, dz = \int_{0}^{1} \int_{0}^{x} 12y^2 \, dy \, dz = \int_{0}^{1} \left[ 4y^3 \right]_{0}^{x} \, dz = \int_{0}^{1} 4x^3 \, dz = 4x^3$.

\textbf{Problem 2: The Missing Factor}
Evaluate the volume of the region bounded by $z = x^2 + y^2$ and $z = 4$.
\textit{Flawed Solution:} 
In cylindrical coordinates, $z$ goes from $r^2$ to $4$, $r$ goes from $0$ to $2$, and $\theta$ from $0$ to $2\pi$. 
$\int_{0}^{2\pi} \int_{0}^{2} \int_{r^2}^{4} 1 \, dz \, dr \, d\theta = \int_{0}^{2\pi} \int_{0}^{2} (4 - r^2) \, dr \, d\theta = \int_{0}^{2\pi} [4r - \frac{r^3}{3}]_{0}^{2} \, d\theta = \int_{0}^{2\pi} (8 - \frac{8}{3}) \, d\theta = 2\pi(\frac{16}{3}) = \frac{32\pi}{3}$.

\textbf{Problem 3: The Even/Odd Trap}
Evaluate $\iiint_E z^3 \, dV$ where $E$ is the unit cube $[0,1] \times [0,1] \times [0,1]$.
\textit{Flawed Solution:} 
The function $z^3$ is an odd function because $f(-z) = -f(z)$. Because we are integrating an odd function over a symmetric region (a perfectly square cube), the integral fundamentally cancels out to zero. Final Answer: $0$.

\textbf{Problem 4: Algebraic Expansion Error}
Evaluate $\int_{0}^{1} \int_{0}^{z} \int_{0}^{y} (y-z)^2 \, dx \, dy \, dz$.
\textit{Flawed Solution:} 
Inner: $\int_{0}^{y} (y^2 - z^2) \, dx = (y^2 - z^2)y = y^3 - yz^2$. 
Middle: $\int_{0}^{z} (y^3 - yz^2) \, dy = [\frac{y^4}{4} - \frac{y^2z^2}{2}]_{0}^{z} = \frac{z^4}{4} - \frac{z^4}{2} = -\frac{z^4}{4}$. 
Outer: $\int_{0}^{1} -\frac{z^4}{4} \, dz = [-\frac{z^5}{20}]_{0}^{1} = -\frac{1}{20}$.

\textbf{Problem 5: Spherical Volume Error}
Evaluate the volume of a sphere of radius 3 using triple integrals.
\textit{Flawed Solution:} 
Using spherical coordinates, the limits are $0 \le \rho \le 3$, $0 \le \phi \le \pi$, and $0 \le \theta \le 2\pi$. 
$V = \int_{0}^{2\pi} \int_{0}^{\pi} \int_{0}^{3} 1 \cdot \rho \, d\rho \, d\phi \, d\theta$. 
Inner: $[\frac{1}{2}\rho^2]_{0}^{3} = \frac{9}{2}$. 
Middle: $\int_{0}^{\pi} \frac{9}{2} \, d\phi = \frac{9\pi}{2}$. 
Outer: $\int_{0}^{2\pi} \frac{9\pi}{2} \, d\theta = 9\pi^2$.

\vspace{0.5cm}
\hrule
\vspace{0.5cm}
\textbf{Answer Key and Corrections:}

\textbf{Correction 1:} 
\textit{Error:} The student swapped the differentials but kept the limits exactly the same; you cannot have $x$ and $y$ as integration limits if you are integrating with respect to them on the inner bounds. 
\textit{Correct Solution:} Keep the original order: Inner $dz \to 24xz$. Middle $dy \to \int_{0}^{x} 24xy \, dy = 12x^3$. Outer $dx \to \int_{0}^{1} 12x^3 \, dx = \left[ 3x^4 \right]_{0}^{1} = 3$.

\textbf{Correction 2:} 
\textit{Error:} The student forgot to multiply the integrand by the required cylindrical Jacobian determinant, $r$.
\textit{Correct Solution:} The integral must be $\int_{0}^{2\pi} \int_{0}^{2} \int_{r^2}^{4} 1 \cdot r \, dz \, dr \, d\theta = \int_{0}^{2\pi} \int_{0}^{2} (4r - r^3) \, dr \, d\theta = \int_{0}^{2\pi} \left[ 2r^2 - \frac{1}{4}r^4 \right]_{0}^{2} \, d\theta = \int_{0}^{2\pi} (8 - 4) \, d\theta = \int_{0}^{2\pi} 4 \, d\theta = 8\pi$.

\textbf{Correction 3:} 
\textit{Error:} The unit cube is on the interval $[0,1]$, which is entirely positive and NOT symmetric about the origin, so the odd-function symmetry rule does not apply.
\textit{Correct Solution:} $\int_{0}^{1} \int_{0}^{1} \int_{0}^{1} z^3 \, dx \, dy \, dz = (1)(1)\left[ \frac{z^4}{4} \right]_{0}^{1} = \frac{1}{4}$.

\textbf{Correction 4:} 
\textit{Error:} The student incorrectly expanded the binomial; $(y-z)^2$ expands to $y^2 - 2yz + z^2$, not $y^2 - z^2$.
\textit{Correct Solution:} Inner: $\int_{0}^{y} (y-z)^2 \, dx = y(y-z)^2 = y^3 - 2y^2z + yz^2$. 
Middle: $\int_{0}^{z} (y^3 - 2y^2z + yz^2) \, dy = \left[ \frac{y^4}{4} - \frac{2y^3z}{3} + \frac{y^2z^2}{2} \right]_{0}^{z} = \frac{z^4}{4} - \frac{2z^4}{3} + \frac{z^4}{2} = \frac{z^4}{12}$.
Outer: $\int_{0}^{1} \frac{z^4}{12} \, dz = \left[ \frac{z^5}{60} \right]_{0}^{1} = \frac{1}{60}$.

\textbf{Correction 5:} 
\textit{Error:} The student used the wrong spherical volume element, writing $\rho$ instead of the correct $\rho^2 \sin\phi$.
\textit{Correct Solution:} $V = \int_{0}^{2\pi} \int_{0}^{\pi} \int_{0}^{3} \rho^2 \sin\phi \, d\rho \, d\phi \, d\theta$. 
Inner: $\left[ \frac{\rho^3}{3} \sin\phi \right]_{0}^{3} = 9\sin\phi$. 
Middle: $\int_{0}^{\pi} 9\sin\phi \, d\phi = \left[ -9\cos\phi \right]_{0}^{\pi} = 18$. 
Outer: $\int_{0}^{2\pi} 18 \, d\theta = 36\pi$.

\end{document}